% !TEX root = ../thesis.tex
\chapter*{PHỤ LỤC}
\addcontentsline{toc}{chapter}{PHỤ LỤC}

\section*{Phụ lục A: Ví dụ kết quả tóm tắt}

Phụ lục này trình bày một ví dụ minh họa về kết quả tóm tắt văn bản tiếng Việt được sinh ra bởi mô hình ViT5 và các thuật toán extractive trên cùng một bài báo từ bộ dữ liệu vietnews, nhằm cung cấp cái nhìn trực quan về chất lượng đầu ra của từng phương pháp.

\textit{Bài báo gốc (trích lược):}\\
\textit{(Sẽ được cập nhật với ví dụ thực tế từ vietnews.)}

\textit{Kết quả tóm tắt của ViT5:}\\
\textit{(Sẽ được cập nhật với kết quả inference thực tế.)}

\textit{Kết quả tóm tắt của các thuật toán extractive:}\\
\textit{(Sẽ được cập nhật sau khi chạy thực nghiệm.)}

\section*{Phụ lục B: Môi trường và cấu hình thực nghiệm}

Bảng dưới đây tóm tắt cấu hình phần cứng và phần mềm được sử dụng trong quá trình thực nghiệm.

\begin{table}[h]
\centering
\renewcommand{\arraystretch}{1.3}
\begin{tabularx}{\textwidth}{l X}
\hline
\textbf{Thành phần} & \textbf{Cấu hình} \\ \hline
GPU & \textit{(sẽ cập nhật)} \\
RAM & \textit{(sẽ cập nhật)} \\
Python & \textit{(sẽ cập nhật)} \\
PyTorch & \textit{(sẽ cập nhật)} \\
Transformers & \textit{(sẽ cập nhật)} \\
Mô hình ViT5 & \texttt{VietAI/vit5-large-vietnews-summarization} \\ \hline
\end{tabularx}
\caption{Cấu hình môi trường thực nghiệm}
\label{tab:env}
\end{table}
